
\documentclass[12pt,a4paper]{article}
\usepackage[utf8]{inputenc}
\usepackage[T1]{fontenc}
\usepackage[brazil]{babel}
\usepackage{amsmath,amssymb}
\usepackage[margin=2.5cm]{geometry}

\begin{document}

\section{Valor Esperado, Variância e Desvio Padrão}

Considere a variável aleatória $X$, que representa o número de requisições que exigem processamento adicional em um lote de 50 requisições independentes.

Conforme estabelecido anteriormente, a probabilidade de uma requisição exigir processamento adicional é $p=0{,}08$. Assim, a variável aleatória segue uma distribuição binomial:

\[
X \sim \operatorname{Binomial}(50,0{,}08).
\]

O objetivo desta etapa é determinar o valor esperado, a variância e o desvio padrão de $X$, utilizando as propriedades da distribuição binomial e verificando os resultados por meio da linguagem R.

\subsection{Valor Esperado}

Para uma variável aleatória binomial, o valor esperado é dado por:

\[
E[X]=np.
\]

Substituindo os parâmetros:

\[
E[X]=50 \cdot 0{,}08.
\]

Portanto,

\[
\boxed{E[X]=4}.
\]

Esse resultado indica que, em média, espera-se que 4 das 50 requisições exijam processamento adicional.

\subsection{Variância}

A variância de uma variável aleatória binomial é definida por:

\[
\operatorname{Var}(X)=np(1-p).
\]

Substituindo os valores:

\[
\operatorname{Var}(X)
=50 \cdot 0{,}08 \cdot (1-0{,}08).
\]

\[
\operatorname{Var}(X)
=50 \cdot 0{,}08 \cdot 0{,}92.
\]

Logo,

\[
\boxed{\operatorname{Var}(X)=3{,}68}.
\]

A variância expressa a dispersão do número de requisições que necessitam de processamento adicional em relação ao valor esperado.

\subsection{Desvio Padrão}

O desvio padrão corresponde à raiz quadrada da variância:

\[
DP(X)=\sqrt{\operatorname{Var}(X)}.
\]

Substituindo:

\[
DP(X)=\sqrt{3{,}68}.
\]

Portanto,

\[
\boxed{DP(X)\approx 1{,}9183}.
\]

O resultado indica uma dispersão típica de aproximadamente 1,92 requisição em torno do valor esperado.

\subsection{Verificação Computacional em R}

Os cálculos foram reproduzidos na linguagem R utilizando as expressões analíticas da distribuição binomial.

A Tabela~\ref{tab:resultados} apresenta a comparação entre os valores calculados manualmente e os resultados computacionais.

\begin{table}[h]
\centering
\begin{tabular}{lcc}
\hline
\textbf{Medida} & \textbf{Manual} & \textbf{R} \\
\hline
Valor esperado & $4$ & $4$ \\
Variância & $3{,}68$ & $3{,}68$ \\
Desvio padrão & $1{,}9183$ & $1{,}9183$ \\
\hline
\end{tabular}
\caption{Comparação dos resultados analíticos e computacionais.}
\label{tab:resultados}
\end{table}

Observa-se concordância entre os resultados teóricos e computacionais, confirmando a aplicação das expressões do valor esperado, da variância e do desvio padrão para a distribuição binomial considerada.

\end{document}
